\section{Introduction}
\label{sec:intro}
\IEEEPARstart{T}{he} Internet of Underwater Things (IoUT) deploys
acoustic multi-hop meshes for environmental observation, resource
exploration, and maritime awareness
\cite{gupta2024uwsn,gola2023taxonomy,saleem2024ai,jiang2023orsurvey}.
Underwater wireless sensor networks (UWSNs) must route under scarce
bandwidth, long delay, and a time-varying multipath channel
\cite{gupta2024uwsn,saleem2024ai,zhu2024trustcomst}.
An increasingly important case is the \emph{mobile} UWSN, in which
autonomous underwater vehicles (AUVs) and unmanned surface vehicles
(USVs) assist monitoring, adaptive sampling, and mobile relay meshes
\cite{gola2023taxonomy,wang2024lgdqn,liu2025mdqn,lu2025wper}.
There the static or moored-relay assumption common in prior routing
work---including LQARP---breaks down, because relative motion
continuously reshapes geometry and instantaneous link quality
\cite{liu2026lqarp,xu2025stgr}.

Routing under mobility is not new, yet two threads remain poorly
aligned with the control-plane cost of keeping link-quality state
accurate while vehicles move.
First, Q-learning and multi-agent reinforcement-learning UWSN
protocols typically assume that neighbor and link state can be
refreshed by periodic beaconing, or they omit an explicit
relative-motion model that would tell a node \emph{when} advertised
link state has become stale
\cite{nandyala2023qtar,zhang2025saqr,liu2024clorP,xu2025stgr,mohanraju2025adrp}.
Periodic HELLO traffic then consumes scarce acoustic bandwidth even
while the channel is temporarily stable.
Second, event-triggered Kalman filtering now schedules filter packets
for IoUT target tracking and remote estimation
\cite{li2026dse,priel2023consensus,sebastian2024eco,deng2023auto},
Kalman SNR prediction is used to adapt a modem
\cite{busacca2024snr}, and event beacons appear in underwater named
data networking \cite{wang2024push,li2024recar}.
Those designs schedule tracking packets, modem updates, or
named-data events.
The subsections below identify which of them, if any, already fire a
routing HELLO from an SNR forecast that is withheld from the
per-tick channel query.
Motivated by these challenges, we review existing
reinforcement-learning underwater routing and event-triggered
estimation, and identify why neither, on its own, resolves that
control-overhead cost.

\subsection{Related Work}
E-LQTP inherits LQARP's multipath channel model, candidate-set
construction, and Bellhop/Aqua-Sim integration unmodified, so this
section does not reopen the simulation-tool argument treated in that
work \cite{liu2026lqarp}.
The subsections that follow isolate the gaps that mobility introduces
on top of that foundation.

\subsubsection{Reinforcement-Learning UWSN Routing}
Nandyala, Kim, and Cho propose QTAR, a Q-learning topology-aware
protocol that selects a next-forwarder from a topology-constrained
candidate set and uses implicit cut-vertex recognition to avoid
energy-wasting detours \cite{nandyala2023qtar}.
Zhang, Chen, and Xie later add service-aware rewards and a holding
time for opportunistic forwarding (SAQR) \cite{zhang2025saqr}.
Yuan \emph{et~al.} nest Q-learning inside unequal clusters
\cite{yuan2023qhier}, and Wang \emph{et~al.} accelerate Q-table
convergence with a depth-aware opportunistic reward
\cite{wang2023rldepth,wang2024qacc}.
Mohanraju and Lokam replace the tabular update with a deep
Q-network that scores residual energy, depth, and node degree
(ADRP-DQL) \cite{mohanraju2025adrp}.
Wang \emph{et~al.} further split the reward by value of information
so that urgent and delay-tolerant packets share the same MARL
forwarder (MARV) \cite{wang2024marv}.
Liu \emph{et~al.} design CLORP, coupling multi-agent reinforcement
learning with opportunistic forwarding and an adaptive learning rate
\cite{liu2024clorP}.
Xu \emph{et~al.} introduce STGR for multi-UUV meshes: a distributed
spatial-temporal graph records neighbor relations, link quality, and
connectivity duration over periodic intervals, and Q-learning then
selects a forwarder on that graph \cite{xu2025stgr}.
He \emph{et~al.} fold a trust score into the same opportunistic
Q-learning loop so that malicious or unstable relays are down-ranked
\cite{he2025trust}.
Li \emph{et~al.} assist the Q-update from an SDN controller
(SDN-QLTR) \cite{li2024sdnqltr}.
LQARP itself folds Bellhop SNR, residual energy, delay, and neighbor
count into a Q-learning reward on a stationary Western Pacific mesh
\cite{liu2026lqarp}.
Its data plane is the parent of E-LQTP; it does not address when a
mobile node must re-advertise that reward state.
Li, Wu, Zhu, and Ren come closest to the control rule studied here.
Their Q-learning link-aware protocol predicts neighbor motion with an
extended Kalman filter and sets the HELLO period to a fraction of the
shortest predicted link expiration time,
$T_H(N_i)=\tau\min_j T_{\mathrm{LET}}(N_i,N_j)$ with
$\tau\in(0,1)$ \cite{li2025qlar}.
High relative motion shortens the period and low relative motion
lengthens it, and the same expiration time enters the reward through
an entropy-weight combination.
That rule already withholds some beacon traffic while a
geometrically predicted link is expected to survive.
It remains a periodic timer scaled by link lifetime, and it spends
the predicted lifetime inside the reward.
It does not compare an SNR forecast, withheld from measurement
updates since the last velocity write, with a channel query taken
every tick.
In the other protocols above, neighbor perception remains a periodic
exchange or a timer-rebuilt graph.

\subsubsection{Mobility-Aware and AUV-Assisted Collection}
Wang \emph{et~al.} plan cooperative multi-AUV collection with
local--global deep Q-learning and a separate emergency multi-hop
route \cite{wang2024lgdqn}.
Liu \emph{et~al.} and Lu and Zuo treat the AUV as a path-planning
agent whose reward includes ocean current and communication capacity
\cite{liu2025mdqn,lu2025wper}.
Nandyala and Cho dispatch an AUV to prevent isolated sub-networks
once isolation is declared \cite{nandyala2025isolated}.
Jiang \emph{et~al.} opportunistically switch between acoustic and
radio hops in a cooperative surface--subsurface mesh
\cite{jiang2023hybrid}.
Yuan \emph{et~al.} jointly optimize multi-objective routes and a
centralized cross-layer schedule of power, links, and next hops
\cite{yuan2024multiobj,yuan2024crosslayer}.
In all of these designs the vehicle is a collector, a repair asset,
or a radio relay---not a source of estimated relative motion that
could gate neighbor broadcasts before advertised SNR goes stale.

\subsubsection{Opportunistic, Clustered, and Information-Centric Forwarding}
Jiang, Han, and Lin survey opportunistic IoUT routing and note that
candidate-set maintenance is still a periodic control cost
\cite{jiang2023orsurvey}.
Zhu \emph{et~al.} cluster relays and predict voids with a Markov
model (RCRP) \cite{zhu2024rcrp}, later adding a GMM--HMM--LSTM
adaptive route \cite{zhu2024gmmhmm}; a companion COMST article
gives trust-management guidelines for the same setting
\cite{zhu2024trustcomst}.
Lee \emph{et~al.} form energy-efficient clusters with a fire-hawk
optimizer and still emit periodic HELLO packets during neighbor
discovery \cite{lee2025fho}.
Wang \emph{et~al.} replace a passive interest with an actively
forwarded beacon when a critical event occurs in underwater named
data networking \cite{wang2024push}.
Li \emph{et~al.} avoid collisions among those NDN paths (RECAR)
\cite{li2024recar}, and Wang, Lu, and Gao cache push and pull
modes for information-centric underwater IoT
\cite{wang2025reliable}.
Zhao \emph{et~al.} add time-reversal opportunistic forwarding
\cite{zhao2026trclor}.
Tian \emph{et~al.} hide source location behind layered candidates
\cite{tian2023lslpr}, and Yuan \emph{et~al.} validate a
source-based disjoint multipath protocol in Zhoushan sea trials
\cite{yuan2024multipath}.
Those beacons and candidate sets are application-level or
path-diversity events, not a routing HELLO tied to a
relative-motion residual.

\subsubsection{Event-Triggered Estimation and Channel Prediction}
Priel and Zelazo develop event-triggered consensus Kalman filtering
for time-varying networks with intermittent observations
\cite{priel2023consensus}.
Sebasti\'{a}n \emph{et~al.} certify an event-triggered distributed
Kalman filter under unknown correlations \cite{sebastian2024eco}.
Deng and Xiong infer the remote state when the trigger decision
itself is hidden \cite{deng2023auto,deng2024lcsys}, and Zhong
\emph{et~al.} schedule multiple sensors over packet-dropping links
\cite{zhong2024tsp}.
Closest in venue and in the use of a predicted state as the trigger
reference, Li~\emph{et~al.} design two media-specific event triggers
for consensus Kalman filtering over a hybrid IoUT
\cite{li2026dse}.
When a surface node does not transmit, neighbors propagate the last
broadcast through the Kalman predict step, yielding a predicted
information pair against which the next innovation is tested.
That construction is the nearest cousin of the forecast track in
Section~III: both ask what a neighbor still believes after a missed
transmission.
Li~\emph{et~al.} reset the predicted information pair when a packet
arrives.
The forecast in Section~III is reset when the local broadcast flag
is high, by copying the neighbor's current odometry, and the copy is
not a demodulated packet.
The scheduled object in \cite{li2026dse} remains a filter packet for
AUV \emph{tracking}, rather than a HELLO that refreshes advertised
link quality for Q-learning forwarding.

A related physical-layer thread predicts SNR itself rather than
scheduling a packet.
Busacca~\emph{et~al.} compare $N$-state Markov and hidden-Markov
models with a scalar Kalman filter on NATO LOON21 traces
\cite{busacca2024snr}.
On those fixed, shallow links the Kalman mean absolute SNR error is
$4.81\,\mathrm{dB}$ on the 11~November 2021 traces,
$3.47\,\mathrm{dB}$ when that model is tested on 14~November, and
$3.24\,\mathrm{dB}$ when it is tested on 18~November.
The $2\,\mathrm{dB}$ cap in Section~III is a design choice of this
paper, placed below those three fixed-link errors.
Busacca~\emph{et~al.} use the filter to adapt a constellation;
they do not prescribe a HELLO threshold.
The dual-track EKF below is a staleness detector for HELLO, not an
SNR-RMSE competitor.

Overall, two issues remain open.
First, the reinforcement-learning and opportunistic protocols
reviewed above, other than QLAR, keep link state current by periodic
beaconing or by a timer-rebuilt graph
\cite{nandyala2023qtar,xu2025stgr,liu2026lqarp,zhu2024rcrp}.
QLAR does test predicted link lifetime, and it uses that lifetime
both to scale a periodic HELLO and to shape the reward
\cite{li2025qlar}.
Second, event-triggered Kalman filtering schedules filter packets
for IoUT tracking and remote estimation
\cite{li2026dse,sebastian2024eco,zhong2024tsp}, Kalman SNR
prediction adapts a modem \cite{busacca2024snr}, and event beacons
appear in UNDN \cite{wang2024push}.
None of those three threads schedules a routing HELLO from the
divergence between an SNR forecast that receives no measurement
update and a channel measurement taken at every tick.

\subsection{Contributions}
The contributions of this article are summarized as follows.

\begin{enumerate}
\item To answer the two gaps jointly, we design a dual-track
relative-motion estimator---a measurement-corrected track and a
forecast track that receives no measurement update---so that an event trigger can be
computed even though the channel-quality measurement is available
every tick.
Its innovation, gated by a persistence condition and a
variance-capped threshold, decides when a node broadcasts the
velocity copy that resets the forecast.
The resulting Event-Triggered Link-Quality and Topology Prediction
protocol (E-LQTP) replaces periodic HELLO traffic with that
motion-aware copy.

\item We investigate whether the same relative-motion estimate can
directly inform the Q-learning reward through a predictive
link-quality term and a drift-aware forwarding term, and we
characterize---rather than tune away---the conditions under which
this does \emph{not} improve routing: a bit-error-rate floor that
creates a scale mismatch with an additive prediction penalty, and a
radial-velocity sign persistence shorter than one hop's transit
time.
This is reported as a scoped, evidence-based negative result.
\end{enumerate}

\subsection{Paper Organization}
The rest of this article is organized as follows.
Section~II introduces the extended system model of a mobile UWSN
under uniform current drift and Rankine-vortex eddy motion.
Section~III details the design of E-LQTP, including the dual-track
extended Kalman filter, the persistence-gated event trigger, and the
scoped Q-learning reward terms.
Section~IV presents simulation results on the Western Pacific site
for both mobile scenarios.
Section~V concludes.
